\chapter{Introducción}
\label{cap:introduccion}

\section{Planteamiento del proyecto}

La exploración robótica de Marte ha mostrado que la distribución y el tamaño de las rocas
en la superficie no son un detalle accesorio, sino una condición que limita lo que una
misión puede hacer. La selección de los sitios de aterrizaje de los \textit{Mars
Exploration Rovers} se apoyó explícitamente en distribuciones tamaño--frecuencia de rocas
derivadas de imágenes orbitales \citep{golombek2003mer}, y el mismo tipo de análisis se
repitió para el \textit{Mars Science Laboratory} \citep{golombek2012msl} y para el módulo
\textit{InSight} \citep{golombek2021insight}. La razón es práctica: una roca de tamaño
suficiente puede impedir el despliegue de un instrumento, obligar a un desvío costoso o
dañar el vehículo. El desgaste prematuro de las ruedas de \textit{Curiosity} sobre terreno
de lecho rocoso con aristas agudas quedó documentado como lección aprendida de la propia
misión \citep{jpl2017wheels}.

Las misiones de superficie han acumulado, a lo largo de dos décadas, cientos de miles de
imágenes tomadas desde el propio terreno. Sobre ese acervo, el conjunto de datos AI4Mars
\citep{swan2021ai4mars, swan2022ai4marsdata} añadió algo poco frecuente: máscaras de
segmentación por píxel para decenas de miles de escenas, producidas mediante anotación
colaborativa y depuradas exigiendo acuerdo entre anotadores. Cada píxel de esas máscaras
lleva una etiqueta de tipo de terreno.

El uso que ha recibido ese material es, casi sin excepción, el de insumo para entrenar
redes neuronales de segmentación. Una revisión sistemática reciente de las técnicas de
segmentación de imágenes marcianas entre 2019 y 2025 \citep{qu2025review} recoge una
sucesión de arquitecturas —redes convolucionales, transformadores, variantes ligeras para
cómputo a bordo— evaluadas todas ellas por su capacidad de reproducir las máscaras. La
máscara funciona como la respuesta correcta contra la cual se mide un modelo.

Este trabajo parte de una observación distinta. Las máscaras de AI4Mars no son solo la
respuesta de un problema de aprendizaje: son en sí mismas un mapa cuantitativo del
terreno. Si un píxel etiquetado como roca informa de que en ese punto de la escena hay
roca, entonces contar píxeles y agrupar los que se tocan permite medir cuánta roca hay y
cómo está organizada, sin entrenar ningún modelo y sin generar máscaras nuevas. Esa
información ya está pagada —el esfuerzo de anotación ya se realizó— y permanece en buena
medida sin explotar.

\section{Pregunta o propósito central}

La pregunta que orienta el trabajo es la siguiente:

\begin{quote}
¿Cómo cuantificar, a partir de las máscaras etiquetadas de AI4Mars, la cobertura de roca
visible y el número aproximado de rocas individuales por imagen, mediante un flujo de
procesamiento de imágenes con parámetros explícitos y resultados reproducibles?
\end{quote}

La pregunta tiene dos mitades de dificultad muy distinta. Medir cobertura es, en esencia,
un conteo de píxeles con un denominador bien elegido. Contar rocas individuales exige
resolver un problema de segmentación de instancias sobre una máscara que no distingue
instancias: cuando dos rocas se tocan, la máscara las registra como una sola región
conectada, y separarlas requiere criterios geométricos. El trabajo aborda las dos y
reporta con igual detalle dónde cada una funciona y dónde no.

\section{Justificación y valor esperado}

El interés de derivar indicadores directamente de las máscaras es triple.

En primer lugar, \textbf{transparencia}. Un indicador calculado con morfología matemática
y división de aguas se puede auditar: cada umbral es un número explícito y cada etapa
intermedia es una imagen que puede inspeccionarse. Un modelo entrenado entrega una
predicción cuya justificación no es directamente legible. Para un uso en el que el
resultado deba defenderse ante un tercero —la planificación de una travesía, por
ejemplo— esa diferencia importa.

En segundo lugar, \textbf{costo}. El procedimiento propuesto corre sobre procesador
convencional, sin unidad gráfica ni conjunto de entrenamiento, y procesa el subconjunto
completo en minutos. Esto lo hace reproducible en cualquier equipo y adecuado como línea
de base contra la cual comparar métodos más costosos.

En tercer lugar, y de manera menos evidente, \textbf{diagnóstico del propio conjunto de
datos}. Al derivar medidas cuantitativas de las máscaras y contrastarlas con las máscaras
de experto que AI4Mars incluye para validación, el procedimiento se convierte en un
instrumento para examinar la calidad de la anotación. El trabajo encontró en esa dirección
uno de sus resultados más relevantes, expuesto en el
Capítulo~\ref{cap:resultados}: la anotación colaborativa sobreestima de forma sistemática
la cobertura de roca. Ese hallazgo tiene consecuencias para cualquier estudio que use
AI4Mars, incluidos los que entrenan modelos.

El valor esperado del trabajo es, por tanto, un procedimiento documentado, un conjunto de
resultados reutilizable con un indicador por imagen y un diagnóstico honesto de los
límites que impone el material de partida.

\section{Objetivos}

\subsection{Objetivo general}

Diseñar, implementar y documentar un procedimiento de procesamiento de imágenes que, a
partir de las máscaras etiquetadas del conjunto de datos AI4Mars, permita cuantificar para
cada imagen la proporción de terreno ocupada por roca visible y el número aproximado de
rocas individuales presentes en las zonas etiquetadas como roca grande, de forma que los
resultados sean reproducibles y comparables entre diferentes escenas.

\subsection{Objetivos específicos}

\begin{description}

  \item[E1.] Estimar, para cada imagen con máscara válida de AI4Mars, la cobertura de
    roca visible, entendida como el porcentaje de píxeles etiquetados de la escena que
    corresponden a roca expuesta o roca grande.

  \item[E2.] Estimar, para cada imagen, el número aproximado de rocas individuales
    presentes en las zonas etiquetadas como roca grande, separando las que aparecen
    pegadas en la máscara mediante componentes conectadas, transformada de distancia y
    división de aguas, y filtrando las regiones resultantes según criterios de tamaño y
    forma.

  \item[E3.] Explorar de forma descriptiva cómo cambian la cobertura y el conteo entre
    distintos subconjuntos y a lo largo del recorrido del rover, en la medida en que esa
    información esté presente en los metadatos disponibles.

  \item[E4.] Implementar el procedimiento en un conjunto de recursos reproducibles que
    incluya código, guiones de ejecución, parámetros empleados, versiones de biblioteca y
    un archivo tabular con los indicadores calculados para cada imagen.

  \item[E5.] Analizar e interpretar los indicadores obtenidos, identificando patrones
    generales, casos atípicos y limitaciones del procedimiento, y plantear líneas de
    trabajo futuras que puedan construirse sobre estos resultados.

\end{description}

\section{Estructura del documento}

El Capítulo~\ref{cap:contexto} caracteriza el contexto del problema, revisa los
antecedentes en abundancia de rocas y en segmentación de terreno marciano, y precisa la
brecha que el trabajo aborda. El Capítulo~\ref{cap:metodos} describe el diseño del
estudio, la fuente de datos, el procedimiento de cálculo con todos sus parámetros y la
estrategia de validación. El Capítulo~\ref{cap:resultados} presenta los resultados de los
cinco objetivos específicos, su validación contra etiquetas de experto y una comparación
con dos enfoques de aprendizaje automático, y discute las implicaciones y limitaciones,
incluido un diagnóstico cuantitativo del origen de los errores del conteo. El
Capítulo~\ref{cap:conclusiones} sintetiza las conclusiones, los aportes y las
recomendaciones. El Anexo~\ref{app:material-complementario} recoge material de soporte:
el esquema del conjunto de resultados, casos adicionales del procedimiento y la
descripción de los recursos reproducibles.
